\section{A paired exclusion tree and reassociation of equal sums}

This section changes only the positive addition producer inside the inherited
retained-total two-stage bit network. It assumes that network's frame compiler,
common flag-basis lemma, copied-centre schedule, and batched recursion contract.
It does not independently establish those inherited algorithmic statements.

\paragraph{Paired leave-one-out recursion.}
Given disjointly supported inputs $z_1,\ldots,z_k$, first remove zero inputs.
If at most two remain, form the total and each omitted sum directly. Otherwise,
rotate the last input to the front, group consecutive inputs in pairs, and form
those pair totals. Recursively compute the total and all omitted pair totals.
For an input in a pair, add the other member of that pair to its omitted pair
total. Undo the rotation. An omitted zero input receives the total.
Induction on $k$ proves that output $i$ is exactly $\sum_{j\ne i}z_j$.
Every addition has disjoint positive supports. This replaces the prefix/suffix
leave-one-out routine in the existing paired pair-exclusion construction;
the global matching order and the retained-point-total construction are unchanged.

\paragraph{Reassociation with exact support equality.}
For every active sum node with support $S$, enumerate existing nodes with
supports $A,B$ satisfying $A\cap B=\varnothing$ and $A\cup B=S$.
Choose another pair only when the integer score described below strictly
increases. Each operand has strictly smaller support cardinality than its sum,
so this cannot introduce a cycle. Prune unused nodes and topologically renumber
by support cardinality. Thus the transformed circuit computes exactly the same
positive sums. It creates no new source support or source subspace.
At the chosen dimension, 368 nodes change their operands, and no active node
is lost or added.

Write $\mathcal I(r)$ for the inherited inner profile: $r$ singleton children
when $2r\le h$, and $h-r$ singleton children together with a child of width
$2r-h$ otherwise. For a node of source-span dimension $d$ with operand dimensions
$a,b$, the changed profile terms are
\[
 \mathcal I(a),\quad\mathcal I(d-a),\quad
 \mathcal I(b),\quad\mathcal I(d-b).
\]
The first and third terms account for the operand fan-out slots; the others
account for the input transitions at the sum node. We maximize the sum of the
squared widths in these four multisets. This is a deterministic integer search
heuristic. No optimality of this score, the resulting producer, or its moment
is claimed; the final complete histogram is certified independently.

\paragraph{Why the inherited source frames apply.}
Every active support consists of triples containing a common point $c$.
Let $t_S$ be the indicator of a triple $S$, and let $G=I-J/9$ be the inherited
Gram form. Then
\[
 \langle t_{\{c,i,j\}},t_{\{c,k,l\}}\rangle_G
 =|\{i,j\}\cap\{k,l\}|.
\]
Hence the Gram matrix on a source family is $B^{\mathsf T}B$, where $B$ is the
unsigned incidence matrix of its pairs on the other $h-1$ points. If a linear
combination is in $\ker B$, its common coordinate also vanishes, since that
coordinate is half the sum of its pair coordinates. Thus the source span is
positive definite and in particular nondegenerate. Its dimension is the
incidence rank, the number of incident vertices minus the number of bipartite
connected components.

Disjoint unions produce nested source spans along every physical slot. Each
side output is still precisely
\[
 \{S:c\in S,\ S\cap(T\setminus\{c\})=\varnothing\},
\]
so its source span lies in $t_T^\perp$. Each retained output remains the full
point star and has dimension $h-1$. The tree replacement has the same positive
common-point property, and reassociation preserves all of its source subspaces.
The nondegenerate nested frame differences therefore remain within the finite
family covered by the inherited generic common flag-basis proposition.

\paragraph{Counts and exact moment certificate.}
At $h=23$ the emitted circuit has
\[
 v=\binom{23}{3}=1771,\qquad m=h^2=529,\qquad N=v^2=3136441,
\]
34,741 addition nodes, 5,313 side outputs, and 23 retained totals. Consequently
\[
 R=40077,\qquad W=2N+2vR=148225616.
\]
The selected auxiliary edges, data-edge profiles, and copied-retained-total
schedule are unchanged. The rank identity is
\[
 s=Wm-N+2vh(h-1)=78410006675,
 \qquad Wm-s=1344189>0.
\]
Let $n_w$ denote the emitted histogram of recursive child widths. The exact
certificate verifies
\[
 \sum_w\frac{w n_w}{Wm}\left(\frac mw\right)^{a_b}<1,
 \qquad a_b=\frac{36943733}{10^{12}}=0.000036943733.
\]
It uses rational upper bounds $U_w\ge\log(m/w)$ and, for $0\le z<4$,
\[
 e^z\le 1+z+\frac{z^2}{2}+\frac{z^3}{6(1-z/4)}.
\]
All arithmetic is rational. The next $10^{-12}$ grid point does not pass this
particular upper-bound certificate; this is not an upper bound on every possible
certificate or construction. The older, looser inequality
$e^z\le(1-z)^{-1}$ with a $10^{-9}$ grid certifies $36941/10^9$ for this same
network. Thus that small numerical difference comes from sharper certification,
whereas the improvement in the histogram comes from the changed producer.

\paragraph{Finite checks and their scope.}
The accompanying checker verifies all 2,047 zero patterns for leave-one-out
lists of lengths zero through ten; independently recomputes every active global
support; checks its common point, every side target and every retained total;
checks all 40,077 nested slot chains; and checks every source rank by the incidence
formula. It additionally checks all 1,970 local active-node ranks by Gaussian
elimination modulo an odd prime. The rational rank and nondegeneracy claims
follow from the incidence argument, rather than from modular experiments alone.
These checks do not certify the full inherited multiplication algorithm.

\section{Assembly choices with a linear phase guard}
The final profiles, including their relaxed strict exponent gaps and $\beta=1/16$,
are specified in \path{notes/latest-assembly.tex} and certified by
\path{research/independent/bit-improvement/assembly_options.py}.
